\loai{Loại 1: Bếp dầu}
\begin{vidu}
	Để đun sôi 4,5 lít nước ở $20^\circ C$ bằng bếp dầu thì ta phải dùng hết $150\ g$ dầu hỏa. Năng suất toả nhiệt của dầu hoả là $q=44.10^{6}\ J/kg$. Biết nhiệt dung riêng của nước là $4200\ J/kg.K$ và khối lượng riêng của nước là 1000 $kg/m^3$. Hiệu suất của bếp dầu hoả này là
	\choice
	{\True $22,9 \%$}
	{$2,29 \%$}
	{$12,9 \%$}
	{$26,9 \%$}
	\loigiai{
		
		Nhiệt lượng nước thu vào $Q_{n}=m_{n} c \Delta t=4,5.4200.(100-20)=1512.10^3\ J$
		
		Nhiệt lượng dầu hỏa tỏa ra $Q=q m_{\text{đ}}=44 . 10^{6}.0,15=6600.10^3\ J$
		
		Hiệu suất $H=\dfrac{Q_{n}}{Q}=\dfrac{1512 . 10^3}{6600 . 10^3} \approx 0,229=22,9 \%$.
		
		
	}
	\haidong{4}\end{vidu}
\begin{vidu}
	Dùng một bếp dầu hoả để đun sôi 2 lít nước từ $15^\circ C$ thì mất 10 phút. Biết rằng chỉ có $40 \%$ nhiệt lượng do dầu toả ra làm nóng nước. Lấy nhiệt dung riêng của nước là $4190\ J/kg. K$, khối lượng riêng của nước là 1000 $kg/m^3$, năng suất toả nhiệt của dầu hoả là $46 . 10^{6}\ J/kg$. Lượng dầu hoả cần dùng cho mỗi phút là
	\choice
	{$0,387\ kg$}
	{$0,0387\ kg$}
	{$0,00738\ kg$}
	{\True $0,00387\ kg$}
	\loigiai{
		
		\begin{itemize}
			\item Nhiệt lượng nước thu vào: $Q_{n}=m_{n} c \Delta t=2.4190.(100-15)=712300\ J$
			\item 		Nhiệt lượng dầu tỏa ra: $Q_{\text{đ}}=\dfrac{Q}{H}=\dfrac{712300}{0,4}=1780750\ J$
			\item 		Khối lượng dầu cần dùng trong 10 phút là: $m_{\text{đ}}=\dfrac{Q_{\text{đ}}}{q}=\dfrac{1780750}{46.10^{6}} \approx 0,0387\ kg$
			\item 		Vậy khối lượng dầu hỏa cần dùng trong mỗi phút là $0,00387\ kg$.
		\end{itemize}
		
	}
	\haidong{8}\end{vidu}
\begin{vidu}
	Để đun sôi một lượng nước bằng bếp dầu có hiệu suất 30\%, phải dùng hết 1 lít dầu. Để đun sôi cũng lượng nước trên với bếp dầu có hiệu suất $20 \%$, thì phải dùng
	\choice
	{2 lít dầu}
	{$\dfrac{1}{2}$ lít dầu}
	{\True 1,5 lít dầu}
	{3 lít dầu}
	\loigiai{
		$H=\dfrac{Q_{n}}{Q_{\text{đ}}}=\dfrac{Q_{n}}{q m}=\dfrac{Q_{n}}{q V \rho} \Rightarrow \dfrac{H_1}{H_2}=\dfrac{V_2}{V_1} \Rightarrow \dfrac{30}{20}=\dfrac{V_2}{1} \Rightarrow V_2=1,5$ lít dầu.
		
	}
	\haidong{6}\end{vidu}

	\loai{Loại 2: Máy phát điện, động cơ điện}
	\begin{vidu} \nguon{1}
		Hiệu suất của một nhà máy nhiệt điện dùng than đá là 0,4. Biết công suất điện của nhà máy là $10\ MW$ và năng suất tỏa nhiệt của than đá là $27 . 10^{6} \ J / kg$. Lấy 1 năm là 365 ngày. Lượng than tiêu thụ hàng năm là
		\choice
		{22,9 nghìn tấn}
		{\True 29,2 nghìn tấn}
		{92,9 nghìn tấn}
		{99,2 nghìn tấn}
		\loigiai{
			$m = \dfrac{\mathcal{P} . t}{H . q} = \dfrac{10^7 . 365 . 24 . 3600}{27 . 10^6 . 0,4} \approx 29200$ nghìn tấn			
		}
	\end{vidu}
	
	\begin{vidu}
		Một nhà máy nhiệt điện tiêu thụ $0,35\ kg$ nhiên liệu cho mỗi kWh điện. Cho biết năng suất tỏa nhiệt của nhiên liệu là $42 . 10^{6}\ J/kg$. Hiệu suất thực của động cơ nhiệt dùng trong nhà máy điện \textbf{gần nhất} với giá trị nào sau đây?
		\choice
		{$38 \%$}
		{$15 \%$}
		{$20 \%$}
		{\True $24 \%$}
		\loigiai{
			\begin{itemize}
				\item $H=\dfrac{A}{Q_1}=\dfrac{\mathcal{\mathcal{P}} t}{q m}=\dfrac{10^3 . 3600}{42 . 10^{6} . 0,35} \approx 0,2449=24,49 \%$.
			\end{itemize}
			
		}
		\haidong{7}\end{vidu}
	\begin{vidu}
		Một máy bơm nước sau khi tiêu thụ hết $8\ kg$ dầu thì đưa được $700\ m^3$ nước lên cao $8\ m$. Biết năng suất toả nhiệt của dầu dùng cho máy bơm này là $4,6.10^{7}\ J/kg$, khối lượng riêng của nước là $1000\ kg/m^3$. Lấy $g=10\ m/s^2$. Hiệu suất của máy bơm là
		\choice
		{\True $15,22 \%$}
		{$24,46 \%$}
		{$1,52 \%$}
		{$2,45 \%$}
		\loigiai{
			
			$A=m g h=V \rho  g h=700.10^3 . 10.8=5,6.10^{7}\ J$
			
			$Q_1=q m=4,6.10^{7}.8=36,8.10^{7}\ J$
			
			$H=\dfrac{A}{Q_1}=\dfrac{5,6 . 10^{7}}{36,8 . 10^{7}} \approx 0,1522=15,22 \%$.
		}
		\haidong{4}\end{vidu}
	
	
	\loai{Loại 3: Động cơ ô tô, xe máy}
	\begin{vidu} \nguon{1}
		Biết hiệu suất của động cơ là $30 \%$, năng suất toả nhiệt của xăng là $4,6 . 10^{7} \ J / kg$, khối lượng riêng của xăng là $700 \ kg / m^3$. Với 4 lít xăng, một xe máy có công suất $1,6 \ kW$ chuyển động với vận tốc $36 \ km / h$ sẽ đi được quãng đường là bao nhiêu km (làm tròn kết quả  đến chữ số hàng đơn vị)?
		\shortans[oly]{ 242} 
		\loigiai{
			\begin{itemize}
				\item TLN: 242
				\item Đổi: $4\ \ell =4.10^{-3} \ m^3$
				\item Khối lượng xăng tiêu thụ là: $m=\rho V=700 . 4 . 10^{-3}=2,8 \ kg$
				\item Nhiệt lượng do xăng tỏa ra là: $Q=mq=2,8 . 4,6 . 10^{7}=12,88 . 10^{7} \ J$
				\item  Công có ích của động cơ: $A=H . Q=0,3 . 12,88 . 10^{7}=38,64 . 10^{6} \ J$
				\item  Thời gian xe máy đã đi là: $t=\dfrac{A}{\mathcal{P}}=24150\ s=6,71\ h$
				\item Quãng đường xe máy đi được: $s=v. t=241,5\ km$.
			\end{itemize}
		}
	\end{vidu}
	
	\begin{vidu} \nguon{1}
		Một ô tô chạy $100 \ km$ với lực kéo không đổi $700 \ N$ thì tiêu thụ hết 6 lít xăng. Biết năng suất toả nhiệt của xăng là $4,6 . 10^{7} \ J / kg$, khối lượng riêng của xăng là $700 \ kg / m^3$. Hiệu suất của động cơ ô tô đó là bao nhiêu \% (làm tròn kết quả  đến chữ số hàng phần mười)?
		\shortans[oly]{36,2 } 
		\loigiai{
			\begin{itemize}
				\item 
				\item Công có ích mà ô tô thực hiện: $A=Fs=700 . 100 . 10^3=7 . 10^{7}\ J$
				\item Nhiệt lượng do 6 lít xăng tỏa ra:
				\item $Q=mq=\rho Vq=700 . 6 . 10^{-3} . 4,6 . 10^{7}=19,32 . 10^{7} \ J$
				\item Hiệu suất của động cơ ô tô đó: $H=\dfrac{A}{Q} . 100 \%=36,23 \%$
			\end{itemize}
		}
	\end{vidu}

\begin{vidu}
	Tính hiệu suất của động cơ ô tô biết rằng khi ô tô chuyển động với vận tốc $72 \ km / h$ thì động cơ có công suất $20 \ kW$ và tiêu thụ 20 lít xăng để chạy $200 \ km$. Cho biết khối lượng riêng của xăng là $700 \ kg / m^3$. Năng suất toả nhiệt của xăng là $4,6 . 10^{7} \ J /kg$.
	\choice
	{\True $31 \%$}
	{$61 \%$}
	{$63 \%$}
	{$36 \%$}
	\loigiai{
		\begin{itemize}
			\item Khối lượng 20 lít xăng: $m=V . \rho=20 . 10^{-3} . 700=14\ kg$
			\item Nhiệt lượng 20 lít xăng toả ra khi bị đốt cháy: $Q_1=q . m=4,6 . 10^{7} . 14=6,44 . 10^{8}\ J$
			\item Thời gian xe chạy 200 km: $t=\dfrac{s}{v}=\dfrac{200}{72}=\dfrac{25}{9} \ h=10000\ s$
			\item Hiệu suất của động cơ: $H=\dfrac{A}{Q_1}=\dfrac{\mathcal{P} . t}{Q_1} \approx 31 \%$
		\end{itemize}
	}
	\haidong{8}\end{vidu}
\loai{Loại 4: Súng bắn}
\begin{vidu} \nguon{2}
	Khi bắn từ súng, viên đạn thép có khối lượng $m=45 \ g$ bắn ra với vận tốc $v=600 \ m/s$. Giả sử rằng $80\%$ năng lượng được giải phóng khi đốt cháy $M=9 \ g$ thuốc súng chuyển thành động năng ban đầu và nhiệt năng của đạn. Biết rằng năng suất tỏa nhiệt của thuốc súng là $q=3\ MJ /kg$, nhiệt dung riêng của thép là $c=500 \ J /kg. K$. Nhiệt độ của đạn tăng lên bao nhiêu độ C?  
	\shortans[oly]{600 }
	\loigiai{
		\begin{itemize}
			\item TLN: 600
			\item Khi thuốc súng cháy, năng lượng được giải phóng là:
			\begin{align*}
				Q_1=q M
			\end{align*}
			Trong đó $M$ là khối lượng thuốc súng
			\item Động năng của đạn là:
			\begin{align*}
				W=\dfrac{m v^2}{2}
			\end{align*}
			\item Sự thay đổi nhiệt năng của đạn là:
			\begin{align*}
				Q_2=c m \Delta t
			\end{align*}
			\item Nếu hiệu suất $H=0,8$, theo điều kiện đề bài ta có:
			\begin{align*}
				H Q_1=W+Q_2
			\end{align*}
			\item Thay số và giải phương trình, ta tìm được:
			\begin{align*}
				\Delta t=\dfrac{0,8 q M-\dfrac{m v^2}{2}}{c. m}=\dfrac{0,8 . 3 . 10^{6} . 9 . 10^{-3}-\dfrac{45 . 10^{-3} . 600^2}{2}}{500 . 45 . 10^{-3}}=600^\circ C
			\end{align*}
		\end{itemize}
	}
\end{vidu}